\fancyhead[L]{\small ADJUSTING THE KNIFE EDGES OF A BALANCE.}
\SourceLeaf{517}\pdfbookmark[0]{Adjusting the Knife Edges of a Balance}{schneiden_waage}
\begin{center}
{\Large [ADJUSTING THE KNIFE EDGES OF A BALANCE.]\par}
\vspace{7mm}\hrule\vspace{2mm}
Göttingische gelehrte Anzeigen. 1837, 13 March.
\vspace{2mm}\hrule\vspace{5mm}
\end{center}
At the meeting of the Royal Society of Sciences on 28 January, Court Councillor \textsc{Gauss} was prompted by Professor \textsc{Weber}'s lecture, reported in issue 22 of this journal, to give a talk on a closely related subject. We give its main points here.

It concerns a new method of adjustment for meeting an essen\-tial requirement in finer beam balances, whose importance so far seems to have been insufficiently appreciated. Such balances have three prismatic knife edges. The one facing down at the centre of the beam rests on a hard horizontal bearing of stone or steel and serves as the axis of rotation as the beam moves. The other two, at the ends of the beam, face up. On each of them rests the support from which the corresponding pan hangs. The supports themselves are made of hardened steel, and their lower surfaces, which rest on the knife edges, are perfectly flat and highly polished.

An essential requirement is that these two outer knife edges be parallel to the central one. Indeed. since the balance is arrested before each exchange of weights in a pan, and the support is thereby lifted off the knife edge, one can never count on the support settling back on the knife edge in \textit{exactly} the same position after the arrest is released. This is harmless if the knife edge concerned is parallel to the central one, but it changes the moment if the knife edges diverge. Imperfect adjustment in this respect is a principal reason why repeated weighings sometimes show considerably greater discrepancies in the results than one
\SourceLeaf{518}would otherwise expect from the excellent workmanship and sensitivity of a balance.

The methods instrument makers have generally used so far to make the knife edges parallel cannot give all the precision one could wish for. Moreover, with fine balances as with astronomical instruments, the finest adjustment is the responsibility of the person who uses the balance. It cannot be demanded of the maker.

The procedure that Court Councillor \textsc{Gauss} has used with the greatest success for this adjustment rests on the following considerations.

As the beam oscillates, the outer knife edge being tested changes its position in space. If this edge is parallel to the (stationary) central one, it remains parallel to itself in all these positions. Things are different if the outer knife edge is out of parallel with the central one. To make the situation easier to picture, let us assume that the outer and central knife edges lie in the same plane, but that their directions diverge away from the observer. In this case, as the beam moves, the outer knife edge will move over a conical surface. Its far end will rise or fall relative to the nearer end as the lever arm on which it lies rises or falls. The same will hold for the support that remains in contact with the knife edge.

It is possible to tell which of these two cases occurs by attaching a plane mirror to the support. The best position for the mirror is with its plane nearly perpendicular to the knife edge, though there is no need to be too particular about this. In the first case the mirror remains parallel to itself as the beam moves. In the second it does not. In the first case, then, the image of an object placed a suitable distance in front of the mirror will remain stationary. In the second it will rise or fall with the lever arm in question, as one readily sees. The reverse would happen if the two knife edges converged away from the observer. The image would then fall when the beam arm rose, and vice versa.

With a very good mirror, even an
\SourceLeaf{519}extremely small shift of the image can be seen reliably and sharply through a telescope. Court Councillor \textsc{Gauss} used as the object a scale graduated in millimetres and set upright about 5 metres in front of the mirror. The telescope, which magnified 35 times, was placed at nearly the same distance. The image of one millimetre thus appeared about 20 seconds of arc across, and tenths of that can still be estimated. As long as the knife edge was short of perfect adjustment, the image of the scale moved up and down quite regularly at the crosshair of the telescope as the beam oscillated.

For readers trained in mathematics it is enough to point out that this method reveals both the direction of the divergence and, with sufficient precision, its magnitude. This information, together with the pitch of the threads of the adjustment screws, makes it possible to carry out the adjustment reliably.

For completeness, a couple of other points may be mentioned here.

When using a fairly large mirror (the one Court Councillor \textsc{Gauss} used was 75 millimetres high and was fastened to the support with a special fixture), the pans must be loaded with sufficiently heavy weights. Otherwise the support would tip sideways.

\Needspace{15\baselineskip}
It was assumed above that the outer knife edge under test lies in the same plane as the central one. Thus if the central edge has been set exactly horizontal and the beam is horizontal, the outer edge is horizontal as well, with at most a sideways divergence. In practice this assumption will generally hold only approximately. In that position the outer edge may be slightly inclined, or one end may be a little higher than the other. In the test described above, this is apparent when the image in the mirror also moves sideways as the beam arm rises, and to the opposite side as it falls. It should be noted, however, that if this error is present in a balance made by a skilled instrument maker, it will in any case be far too small to cause any appreciable error in the results of the weighings. This is why even the best balances have no means of adjustment to eliminate this component of the departure from parallelism.

